\documentclass[10pt,a4paper]{amsart}
\usepackage[margin=19mm]{geometry}
\usepackage{amsmath,amssymb,mathtools}
\usepackage[scr=rsfs,cal=euler]{mathalfa}
\usepackage{xfrac}
\usepackage[unicode,hidelinks,bookmarks=false]{hyperref}
\newcommand{\ie}{i.e.\ }
\newcommand{\cf}{cf.\ }
\newcommand{\resp}{respectively\ }
\newcommand{\Subsection}{\subsection}
\providecommand{\nobreakdash}{\nobreak\hbox{-}}
\setlength{\emergencystretch}{2em}
\multlinegap0pt
\usepackage{bm}
\usepackage{bbm}
\usepackage{dsfont}
\usepackage{xspace}
\usepackage{cancel}
\usepackage{mathtools}
\usepackage{amsthm}
\makeatletter
\@ifundefined{theorem}{\newtheorem{theorem}{Theorem}}{}
\@ifundefined{proposition}{\newtheorem{proposition}[theorem]{Proposition}}{}
\makeatother
\DeclareMathOperator{\HF}{HF}
\DeclareMathOperator{\ann}{Ann}
\newcommand{\K}{\mathbb{K}}
\title[Artinian Gorenstein algebras with Macaulay dual generator wi]{Artinian Gorenstein algebras with Macaulay dual generator with fixed Waring rank}
\author{Janaine Martins, Rosa M. Mir\'o-Roig}
\date{Source: arXiv:2608.12621v1 (CC BY 4.0)}
\begin{document}
\maketitle

\noindent\emph{Source and licence.} Artinian Gorenstein algebras with Macaulay dual generator with fixed Waring rank. Version arXiv:2608.12621v1 (CC BY 4.0), distributed under the Creative Commons Attribution 4.0 International licence, as recorded in the arXiv licence metadata for this exact version and shown on the abstract page https://arxiv.org/abs/2608.12621v1. Full rights notice: licenses/arxiv-2608.12621-CC-BY-4.0.txt.\par\smallskip

\noindent\textit{Editorial scope.} Counted unit: the Proposition whose source label is \texttt{h\_vector} in the article (source0.tex lines 515--527, source label \texttt{h\_vector}), delivered together with the article's own complete original proof of it (source0.tex lines 528--568, one \texttt{proof} environment of 41 lines). No mathematics is written, repaired or re-derived: statement and proof are the article's own text line for line. Required context retained uncounted: computingHilbertFunctionGorenstein (lines 326--335). Every definition, display and label of those lines is retained unchanged. Omitted from this package and not redistributed: 1--325, 336--514, 569--1086. Nothing inside a retained excerpt is omitted or altered: every retained body is delivered exactly as the article's lines are written, so no omission receipt of any kind accompanies this package. Citations: the retained excerpts carry no citation command at all, so no bibliography entry of the article is transcribed and no reference is redistributed. Reference adaptation: none was needed and none was made; every reference inside the delivered unit resolves inside it, so no unresolved reference or citation is produced. Printed slips kept literal: no printed word, symbol, hypothesis or number of the counted statement or of its proof was altered. Layout and class: the article's own class is \texttt{amsart} with packages including amssymb, color, tikz, inputenc, fontenc, tikz-cd, lipsum, adjustbox, graphicx, comment, xcolor, centernot, mathtools, stmaryrd; the wrapper is the lane's pinned \texttt{amsart} editorial wrapper at 10pt on A4, which restates the theorem-like environments the retained text needs and adds the article's own macro definitions that the retained text needs (\texttt{\textbackslash HF}, \texttt{\textbackslash K}, \texttt{\textbackslash ann}), with extra wrapper packages (bm, bbm, dsfont, xspace, cancel, mathtools). The delivered numbering is the unit's own. Assets: no retained span contains a figure environment, a graphics-inclusion command, a PDF-page inclusion, a table or a TikZ picture; no image file is copied into this package, the package declares no asset dependency and no image file is redistributed with it. Licence: version arXiv:2608.12621v1 of this article is distributed under the Creative Commons Attribution 4.0 International licence, as recorded in the arXiv licence metadata for this exact version and shown on the abstract page https://arxiv.org/abs/2608.12621v1. A full rights notice is delivered at licenses/arxiv-2608.12621-CC-BY-4.0.txt. Characters: every retained excerpt is pure ASCII, so the delivered engine, XeLaTeX with its default Latin Modern text font, reports no missing character.
\begin{equation}
    \label{computingHilbertFunctionGorenstein}
\begin{array}{lcl}
    \HF_A(t) &= &\dim_{\K}[A]_t \\
    &= &\dim_{\K}[A]_{c-t} \\
    & =&\dim_{\K}[I^{-1}]_{c-t}  \\
    &= &\dim_{\K} \left\langle \frac{\partial ^{t}F}{\partial X_1^{t_1}\cdots \partial X_n^{t_n}} \mid t_1+\cdots +t_n= t
    \right\rangle .
\end{array}
\end{equation}

\begin{proposition}\label{h_vector}We fix integers $d$, $s\ge 1$ and $n\ge 2$ such that $s\le {n+d-1\choose n-1}$. Let $A_{F_{s}} = R/\ann(F_{s}),$ be the AG $\K$-algebra  of codimension $n$, socle degree $d$ and Macaulay dual generator 
\begin{eqnarray*}
F_{s}=\ell_1^d+\cdots+\ell_s^d
\end{eqnarray*}
with $\ell_1,\dots,\ell_s$ general linear forms in $[S]_1$.  
Then the Hilbert function of $A_{F_{s}}$ is given by
$$
h_k=\dim_{\mathbb K}[A_{F_{s}}]_k
=\begin{cases}\min\left\{\binom{n+k-1}{k},\, s\right\} \text{ for }  0\leq k\leq\lfloor\frac{d}{2}\rfloor \\
\text{symmetry}  \text{ for } \lfloor\frac{d}{2}\rfloor < k \leq d  .
\end{cases}
$$
\end{proposition}

\begin{proof}
  By the Equation (\ref{computingHilbertFunctionGorenstein}), the $k$-th homogeneous component $[A_{F_s}]_k$ of  $A_{F_s}$ can be  identified with the vector subspace of $S_{d-k}$
generated by the  $k$-th  partial derivatives of $F_s$. On the other hand, for a linear general form
\begin{eqnarray*}
\ell_i=a_{i}^1X_1+\cdots+a_{i}^nX_n,
\end{eqnarray*}
we obtain:
\begin{eqnarray*}
\frac{\partial^k}{\partial x_{j_1}\cdots\partial x_{j_k}}(\ell_i^d)
=
\frac{d!}{(d-k)!}
a_{i}^{j_1}\cdots a_{i}^{j_k}\ell_i^{d-k},
\end{eqnarray*}
and, we get:

$$
\begin{array}{rcl}
\frac{\partial^kF_s}{\partial x_{j_1}\cdots\partial x_{j_k}} &
= & 
\frac{\partial^k}{\partial x_{j_1}\cdots\partial x_{j_k}}(\sum _{i=1}^s\ell_i^d )\\  \\
& = &
\frac{d!}{(d-k)!}\sum _{i=1}^s
a_{i}^{j_1}\cdots a_{i}^{j_k}\ell_i^{d-k}.
\end{array}
$$

Moreover, we have
$$[A_{F_s}]_k=\left \langle  \frac{\partial^kF_s}{\partial x_{j_1}\cdots\partial x_{j_k}} \ \mid  \ 1\le j_1\le \cdots \le j_k\le n   \right\rangle .$$
Since $\left[A_{F_s}\right]_k=\operatorname{Span}\left\{\ell_1^{d-k}, \ldots, \ell_s^{d-k}\right\},$ its dimension is bounded above by both $\operatorname{dim}_{\mathbb{K}} R_k=\binom{n+k-1}{k}$ and by the number of generators $s$. Hence, $h_k \leq \min \left\{\binom{n+k-1}{k}, s\right\}.$

Now suppose first that
$s \leq\binom{ n+k-1}{k}$. As the linear forms $\ell_1, \ldots, \ell_s$ are general, the powers $\ell_1^{d-k}, \ldots, \ell_s^{d-k}$
are linearly independent in $S_{d-k}$. Therefore, $h_k=s.$

On the other hand, if $s > \binom{n+k-1}{k},$ then the only possible linear relations among the $k$-th derivatives are those coming from the space of differential operators itself. Since the linear forms are general, no additional relations occur, and consequently $h_k=\operatorname{dim}_{\mathbb{K}} R_k=\binom{n+k-1}{k}.$

Thus, for any $k$ such that $0\leq k\leq\lfloor\frac{d}{2}\rfloor$, it holds that
$$h_k=\dim_{\mathbb K}[A_{F_{s}}]_k
=\min\left\{\binom{n+k-1}{k},\, s\right\} $$
and, since $A_{F_s}$ is an AG $\K$-algebra  of socle degree $d$, the Hilbert function of $A_{F_s}$ is symmetric (i.e. $h_k = h_{d-k}$) and this finishes the proof. 
\end{proof}

\end{document}
